\chapter{Resultados auxiliares sobre espacios métricos}

\section{Convergencia y continuidad uniformes}

\transclude{convergencia-uniforme}

\transclude{continuidad-uniforme}

\begin{prop}\label{prop:convergencia-uniforme-continuidad-uniforme}
    Sean $X, Y$ \exhyperref[defn:esp-metrico]{esp-metrico}{espacios métricos}, $f \colon X \to Y$ y sea $\left(f_n\right)_{n \in \N}$ una sucesión de funciones $f_n \colon X \to Y$ \exhyperref[defn:continuidad-uniforme]{continuidad-uniforme}{uniformemente continuas} que \exhyperref[defn:convergencia-uniforme]{convergencia-uniforme}{converge uniformemente} a $f$. Entonces, $f$ es uniformemente continua.
\end{prop}
\begin{dem}
    Sea $\varepsilon > 0$. Entonces, $\exists N \in \mathbb{N} : \forall x \in X : d_Y(f_N(x), f(x)) < \frac{\varepsilon}{3}$.

    Por otro lado, como $f_N$ es uniformemente continua, $\exists \delta > 0 : \forall x_1, x_2 \in X : d_X(x_1, x_2) < \delta \implies d_Y(f_N(x_1), f_N(x_2)) < \frac{\varepsilon}{3}$. Por tanto, si $d_X(x_1, x_2) < \delta$, por la \exhyperref[defn:metrica:desigualdad-triangular]{metrica}{desigualdad triangular}, se tiene que
    \begin{align*}
        d_Y(f(x_1), f(x_2)) &\leq d_Y(f(x_1), f_N(x_1)) + d_Y(f_N(x_1), f_N(x_2)) + d_Y(f_N(x_2), f(x_2)) \\
        &< \frac{\varepsilon}{3} + \frac{\varepsilon}{3} + \frac{\varepsilon}{3} = \varepsilon.
    \end{align*}
    Luego $f$ es uniformemente continua.
\end{dem}

\subsection{Convergencia uniforme sobre compactos}

\transclude{convergencia-uniforme-compactos}

\section{Combinaciones y envolventes convexas}

\begin{defn}[Combinación convexa]\label{defn:combinacion-convexa}
    Sea $V$ un $\R$-\exhyperref[defn:esp-vectorial]{esp-vectorial}{espacio vectorial} y $\left\{x_i\right\}_{i=1}^n \subset V$, $y \in V$ es una combinación convexa de $\left\{x_i\right\}_{i=1}^n$
    \[
    \iff \exists \left\{\lambda_i\right\}_{i=1}^n \subset \mathbb{R}_{\geq 0} : \sum_{i=1}^n \lambda_i = 1 \we y = \sum_{i=1}^n \lambda_i x_i
    \]

    Sea $A \subset V$, $y \in V$ es una combinación convexa de elementos de $A \iff \exists \left\{a_i\right\}_{i=1}^n \subset A$ tal que $y \tex{ es una combinación convexa de } \left\{a_i\right\}_{i=1}^n$.
\end{defn} %ORDENAR enviar al anexo

\begin{defn}[Envolvente convexa]\label{defn:envolvente-convexa}
    Sea $V$ un $\R$-\exhyperref[defn:esp-vectorial]{esp-vectorial}{espacio vectorial} y $A \subset V$, $\operatorname{conv}(A)$ es la envolvente convexa de $A \iff \operatorname{conv}(A)$ es el \exhyperref[defn:con-convexo]{con-convexo}{convexo} más pequeño que contiene a $A$.
\end{defn} %ORDENAR enviar al anexo

\begin{prop}\label{prop:envolvente-convexa-combinacion-convexa}
    Sea $V$ un $\R$-\exhyperref[defn:esp-vectorial]{esp-vectorial}{espacio vectorial} y $A \subset V$. Entonces,
    \[
    \operatorname{conv}(A) = \left\{y \in V : v \tex{ es una \exhyperref[defn:combinacion-convexa]{combinacion-convexa}{combinación convexa} de elementos de } A \right\}.
    \]
\end{prop}%ORDENAR enviar al anexo


%ORDENAR: reposicionar este teorema
\begin{teo}[Stone--Weierstrass en $\mathbb{T}$]\label{teo:stone-weierstrass-toro}
    Sean $\eta \in \mathcal{C}(\mathbb{T})$ y $\varepsilon > 0$. Entonces, existen $N \in \N$ y $\left(a_k\right)_{k=-N}^N \subset \C$ tales que
    \[
    \max_{\xi \in \mathbb{T}} \abs{\eta(\xi) - \sum_{k=-N}^N a_k \xi^k} < \varepsilon.
    \]
\end{teo}